% !TeX root = bash.tex

\section{IV. Regex}
\begin{frame}
\frametitle{What's a regex}
\textbf{Regex}: ``regular expressions''
\newline \newline
It is a string… to match patterns in other strings.
\end{frame}

\begin{frame}[fragile]
\frametitle{What can (and can't) a regex do}
A regex can (attempt to):
\begin{itemize}
    \item Match a cell phone number: \cverb|[0-9]{11}|
    \item Match a .com domain: \cverb|([A-Za-z0-9-]+\.)+com|
    \item Match \textit{Star Wars} subtitles but not \textit{Star Trek}:
        \cverb!m | [tn]|b!\footnote{Credit: \url{https://xkcd.com/1313/}}
	\footnote{Explained: \url{https://www.explainxkcd.com/wiki/index.php/1313:_Regex_Golf}}
        \footnote{The art of regex golf: \url{https://alf.nu/RegexGolf}}
\end{itemize}
It cannnot:
\begin{itemize}
    \item Parse Python code
    \item Detect homework plagiarism
    \item Moderate a Minecraft server
\end{itemize}
\end{frame}

\begin{frame}
\frametitle{Regex is a mess}
\begin{block}{Quote}
I define UNIX as 30 definitions of regular expressions living under one roof.
\begin{flushright}
    — Donald Knuth\footnote{Digital Typography, ch. 33, p. 649 (1999)}
\end{flushright}
\end{block}

Two dominant standards:
\begin{itemize}
    \item \textbf{ERE} (Extended RegEx, POSIX compliant)
    \item \textbf{PCRE} (Perl Compatible RegEx, used in Perl and Python)
\end{itemize}
Our focus today will be ERE.
\end{frame}

\begin{frame}[fragile]
\frametitle{Regex patterns: \tt{.}}
The dot is the simplest pattern:
\begin{table}
    \centering
    \begin{tabular}{ll}
        \textbf{Pattern}    & \textbf{Matches} \\
        \cverb|.|            & any character \\
    \end{tabular}
\end{table}

\begin{example}
    \cverb|c.t| matches \cverb|cat|, \cverb|cut|, \cverb|c t|, etc.
\end{example}

\begin{block}{Note}
    \cverb|\.| matches a literal dot. Same goes for \cverb|\(|, \cverb|\[|, etc.
\end{block}
\end{frame}

\begin{frame}[fragile]
\frametitle{Regex patterns: \tt{[]}}
Brackets match any of the characters inside, but \cverb|^| and \cverb|-| are special:
\begin{table}
    \centering
    \begin{tabular}{ll}
        \cverb|[aeiou]|      & any vowel \\
        \cverb|[^aeiou]|     & anything but a vowel \\
        \cverb|[0-9]|        & any digit \\
        \cverb|[^A-Za-z]|    & anything but letters \\
    \end{tabular}
\end{table}

\begin{example}
    \begin{itemize}
        \item \cverb|VV1[58]6| matches \cverb|VV156| and \cverb|VV186|
        \item \cverb|[A-C][01][0-9]| matches \cverb|A00| up to \cverb|C19|
    \end{itemize}
\end{example}

\begin{block}{Note}
    To escape dashes, anchors, and brackets, use \cverb|\-|, \cverb|\^|, \cverb|\[|, and \cverb|\]|, respectively.
\end{block}
\end{frame}

\begin{frame}[fragile]
\frametitle{Regex patterns: character classes}
\begin{table}
    \centering
    \begin{tabular}{ll}
        \cverb|\w|           & \cverb|[A-Za-z0-9_]| \\
        \cverb|\W|           & anything \cverb|\w| does not match \\
        \cverb|\s|           & whitespace (space, tab, linebreak, etc) \\
        \cverb|\S|           & anything but whitespace \\
    \end{tabular}
\end{table}

\begin{block}{Note}
    Character classes in brackets like \cverb|[\w\s]| won't work.
\end{block}
\end{frame}

\begin{frame}[fragile]
\frametitle{Regex patterns: \tt{| ()}}
Vertical bars separate patterns, and matches one of them.
Parentheses can be used to group patterns.
\begin{table}
    \centering
    \begin{tabular}{ll}
        \cverb!stan|kyle|eric|kenny!        & \cverb|stan|, \cverb|kyle|, \cverb|eric|, or \cverb|kenny| \\
        \cverb!(ls|cd|rm -r) dir!    & \cverb|ls dir|, \cverb|cd dir|, or \cverb|rm -r dir| \\
        \cverb!ls|cd|rm -r dir!      & \cverb|ls|, \cverb|cd|, or \cverb|rm -r dir| \\
    \end{tabular}
\end{table}
\end{frame}

\begin{frame}[fragile]
\frametitle{Regex patterns: repeat}
A repeated pattern can be matched:
\begin{table}
    \centering
    \begin{tabular}{ll}
        \cverb|A?|           & zero or one A \\
        \cverb|A+|           & one or more A's \\  % Positive closure
        \cverb|A*|           & zero or more A's \\ % Kleene closure
        \cverb|A{6}|         & 6 A's \\
        \cverb|A{4,6}|       & 4 to 6 A's \\
        \cverb|A{4,}|        & at least 4 A's \\
    \end{tabular}
\end{table}

\begin{example}
    \begin{itemize}
        \item \cverb|hh.*| matches \cverb|hh|, \cverb|hhgg|, \cverb|hhjj|, \cverb|hh!!!!!!|, etc.
        \item \cverb|a(ba)+| matches \cverb|aba|, \cverb|ababa|, \cverb|abababababa|, etc.
        \item \cverb|[0-9]{1,3}(,[0-9]{3})*| matches \cverb|6|, \cverb|114,514|, \cverb|1,919,810|, etc.
    \end{itemize}
\end{example}
\end{frame}

\begin{frame}[fragile]
\frametitle{Regex patterns: location}
These do not match literal characters. Instead, they specify the location of
the character before or after it.
\begin{table}
    \centering
    \begin{tabular}{ll}
        \cverb|^|            & beginning of string \\
        \cverb|$|            & end of string \footnote{Beginning or end of line sometimes} \\
        \cverb|\b|           & word boundary (boundary between \cverb|\w| and \cverb|\W|) \\
        \cverb|\B|           & not word boundary \\
    \end{tabular}
\end{table}

\begin{example}
    \begin{itemize}
        \item \cverb|^$| matches an empty string only
        \item \cverb|vsky$| matches russian names
        \item \cverb|.*coke\b.*| matches \cverb|would you like some coke?|, \cverb|yes i'd love to have some coke|
            but not \cverb|i drink twelve cokes a day|
    \end{itemize}
\end{example}
\end{frame}

\begin{frame}[fragile]
\frametitle{Quiz: Does it match?}
What strings does this regex match? \newline

\Large \cverb!^cat|cat$! \normalsize

\begin{itemize}
    \item \cverb|cat|                      % yes
    \item \cverb|^cat$|                    % no
    \item \cverb|cats|                     % yes
    \item \cverb|cat /home/anon/30_days_to_building_| \\
        \cverb|a_healthy_relationship_with_your_cat|           % yes
    \item \cverb|I have a cat.|            % no
    \item \cverb|Cats are the best.|       % no
    \item \cverb|Concatenate these files|  % no
\end{itemize}
\end{frame}

\begin{frame}[fragile]
\frametitle{Quiz: Does it match?}
What matches \cverb|021 54749110| but not \cverb|02154749110|?
\begin{itemize}
    \item \cverb|\021\s+[0-9]{8}| % this one
    \item \cverb|\021\s*[0-9]{8}|
\end{itemize}
\end{frame}

\begin{frame}[fragile]
\frametitle{Quiz: Does it match?}
Construct a regex that matches \cverb|I'm from SJTU|, \cverb|I'm from SJTUJI|, and \cverb|I'm from SJTUGC|.
% I'm from SJTU(JI|GC)?
\end{frame}

\begin{frame}[fragile]
\frametitle{But how to use a regex, anyway?}
Try this in \tt{04-regex/}:
\footnote{For Mac users, instead of \ctexttt{grep} you might have to type \ctexttt{ggrep}}
\begin{lstlisting}[language=bash]
$ grep -E '.+\..+@sjtu.edu.cn' faculty
\end{lstlisting}
\pause
\begin{block}{Observation}
    The regex matches all email addresses in the file \tt{faculty}
    that look like \cverb|firstname.lastname@sjtu.edu.cn|.
    \newline \newline
    \cverb|-E| stands for Extended regex.
\end{block}
\end{frame}

\begin{frame}[fragile]
\frametitle{One more example}
Try this in \tt{04-regex/}:
\begin{lstlisting}[language=bash]
$ grep -oE '^[^@]{,8}' faculty
\end{lstlisting}
\pause
\begin{block}{Observation}
    \cverb|@sjtu.edu.cn| are all gone, and each line is at most 8 characters long.
\end{block}
\begin{block}{Explanation}
    \begin{tabular}{ll}
        \cverb|-o|   & Print only the matched parts \\
        \cverb|^|    & From beginning of each line \\
        \cverb|[^@]| & Keep any character except @ \\
        \cverb|{,8}| & Until we reach length 8
    \end{tabular}
\end{block}
\end{frame}

\begin{frame}[fragile]
\frametitle{Your turn}
Extract all course codes from \tt{04-regex/courses}. \newline

\begin{example}
    \begin{tabular}{lcl}
        \cverb|VG100 Introduction to Engineering| & & \cverb|VG100| \\
        \cverb|VM020 Machineshop Training| & $\Longrightarrow$ & \cverb|VM020| \\
        \cverb|VP140 Physics I| & & \cverb|VP140|
    \end{tabular}
\end{example}
\pause
\begin{block}{Solution (naïve version)}
\begin{lstlisting}[language=bash]
$ grep -oE 'V[A-Z][0-9]+' courses
\end{lstlisting}
\end{block}
\end{frame}

\begin{frame}[fragile]
\frametitle{Find and replace with \tt{sed}}
\cverb|sed| is a powerful tool for transforming text.
\footnote{For Mac users, this may be called \tt{gsed}.}
We will be using one very specific syntax for substitution:\footnote{
    When the pattern/replacement contains slashes, you can use things like
    \tt{!} and \tt{,} as delimiters.
}
\begin{lstlisting}[language=bash]
$ COMMAND | sed -E 's/FIND/REPLACE/FLAGS'
$ sed -E 's/FIND/REPLACE/FLAGS' FILE
\end{lstlisting}
\cverb|-E| again stands for Extended regex, and flags are optional.
This command redacts all the IPv4 addresses in the file \tt{ipv4}:
\begin{lstlisting}[language=bash]
$ sed -E 's/([0-9]{1,3}\.){3}[0-9]{1,3}/redacted/g' \
    ipv4
\end{lstlisting}
\begin{block}{Observe}
    What will happen without the \cverb|g| at the end?
\end{block}
\end{frame}

\begin{frame}[fragile]
\frametitle{The \texttt{/g} flag}
\cverb|sed| by default changes only the first occurrence on each line.
To change every occurrence, use the \cverb|/g| flag for global replacement.
\end{frame}

\begin{frame}[fragile]
\frametitle{Capturing groups}
A \textbf{capturing group}, or simply \textbf{group}, is a pattern inside
parentheses that mark a region of text you want to keep in the replaced text.
\newline \newline
It can be accessed with a respective \textbf{backreference} which looks like
\cverb|\1|, \cverb|\2|, up to \cverb|\9|.

When nested, the position of \cverb|(| determines order, so the outside group is
\cverb|\1| and the inside is \cverb|\2|.
\end{frame}

\begin{frame}[fragile]
\frametitle{Capturing groups with \tt{sed}}
What if you only want to redact the subnet (i.e. last part) of the IP addresses?
\begin{lstlisting}[language=bash]
$ sed -E 's/(([0-9]{1,3}\.){3})[0-9]{1,3}/\1xxx/g' \
    ipv4
\end{lstlisting}
\begin{block}{Observation}
    IP addresses like \cverb|192.168.1.1| become \cverb|192.168.1.xxx|
\end{block}
\end{frame}

\begin{frame}[fragile]
\frametitle{Your turn}
From \tt{04-regex/courses}, select 100- and 200-level math courses and convert
legacy ``VV'' course codes into modern ``MATH'' codes.
Do not print other courses.
\begin{example}
    \begin{tabular}{lcl}
        \cverb|VV156| & & \cverb|MATH1560J| \\
        \cverb|VV214| & $\Longrightarrow$ & \cverb|MATH2140J| \\
        \cverb|VV417| & & \textit{not printed}
    \end{tabular}
\end{example}
\pause
\begin{block}{Solution}
\begin{lstlisting}[language=bash]
$ grep -E 'VV[12]' courses | \
    sed -E 's/VV([0-9]{3})/MATH\10J/'
\end{lstlisting}
\end{block}
\end{frame}

\begin{frame}
\frametitle{When not to use regex?}
\begin{columns}
    \begin{column}{0.7\textwidth}
        \includegraphics[width=\textwidth]{regex_explained}
    \end{column}
    \begin{column}{0.3\textwidth}
        If your regex:
        \begin{itemize}
            \item is 50 characters long
            \item handles lots of Unicode
            \item has too many backslashes
        \end{itemize}
        consider something else.
    \end{column}
\end{columns}
\end{frame}

\begin{frame}[fragile]
\frametitle{Atrocities under regex's name}
If you match a Chinese resident ID with \cverb|[0-9]{18}|, people whose
ID ends with X will be mad at you.

Apart from ignorance, people also abuse regex for fun. This is a regex
that matches a regex:
\begin{lstlisting}
/\/((?![*+?])(?:[^\r\n\[/\\]|\\.|\[(?:[^\r\n\]\\]|\\.)*\])+)\/((?:g(?:im?|mi?)?|i(?:gm?|mg?)?|m(?:gi?|ig?)?)?)/
\end{lstlisting}

And here's one that matches integers that are divisible by 3:
\begin{lstlisting}
^(?:[0369] |
(?:[147](?:[147][0369]*[258]|[0369])*[258]) |
(?:[258](?:[258][0369]*[147]|[0369])*[147]))+$
\end{lstlisting}

(These two examples are not ERE.)
\end{frame}

\begin{frame}
\frametitle{Beyond the workshop}
Regex101 (\url{https://regex101.com/}) is an online regular expression
evaluator that supports ERE, PCRE, and many others.
\end{frame}
